%auto-ignore

\usepackage{fancyvrb,fvextra}
\usepackage{xcolor} % For coloring text

\usepackage[protrusion=false]{microtype}
\usepackage{graphicx}
\usepackage{subfigure}
\usepackage{booktabs} 
\usepackage{color}
\usepackage{algpseudocode}
\usepackage{algorithm}

\usepackage{natbib}
\usepackage{hyperref}
\usepackage{ulem}
\usepackage{amsmath}
\usepackage{amssymb}  
\usepackage{mathtools}
\usepackage{amsthm}
\usepackage{cleveref}
\usepackage{bm}
\usepackage{tcolorbox}
\usepackage{listings}
\usepackage{setspace}
\usepackage{tikz}


\usepackage{listings}
\lstset{
    basicstyle=\small\ttfamily,
    breaklines=true,
    breakatwhitespace=true,
    columns=flexible,
    keepspaces=true,
    xleftmargin=0.3cm,
    xrightmargin=0.3cm,
    breakindent=0pt,
    escapeinside={(*@}{@*)}
}

\usetikzlibrary{positioning, fit, arrows.meta}
%  \doublespace
\DefineVerbatimEnvironment{MyVerbatim}{Verbatim}{
  commandchars=@\{\},
  breaklines=true
}


% \lstset{
%     % basicstyle=\ttfamily\small,     % Set font to monospaced and small
%     breaklines=true,                % Enable automatic line breaking
%     % frame=single,                   % Add a single-line frame around the code
%     % numbers=left,                   % Show line numbers on the left
%     moredelim={[is][keywordstyle]{@@}{@@}},
%     % numberstyle=\tiny\color{gray},  % Style for line numbers
%     % xleftmargin=15pt,               % Add left margin
%     % xrightmargin=15pt,              % Add right margin
%     showstringspaces=false         % Don't show spaces in strings
% }

\lstdefinestyle{customstyle}{
    moredelim={[is][keywordstyle]{@@}{@@}},  % Set the delimiters and styling
    keywordstyle=\color{blue}\textbf,               % Example keyword style
    breaklines=true,  
    basicstyle=\ttfamily
}

% Apply the custom style globally
\lstset{style=customstyle}
\lstset{
  literate={```}{{\textasciigrave\textasciigrave\textasciigrave}}2,
}

\newtcolorbox{mybox}[1][]{
    title=#1,
    fonttitle=\small,
    fontupper=\small,
    left=2mm,
    right=2mm,
    top=1mm,
    bottom=0mm,
}

\newtheorem{assumption}{Assumption}
\newtheorem{definition}{Definition}
\newtheorem{theorem}{Theorem}
\newtheorem{lemma}[theorem]{Lemma}
\newtheorem{observation}{Observation}
\crefname{observation}{Observation}{Observations}

\DeclareMathOperator*{\E}{\mathbb{E}}

\DeclareMathOperator*{\argmax}{\mathrm{argmax}}
\DeclareMathOperator*{\argmin}{\mathrm{argmin}}


\def\midd{\middle\|}

\newcommand{\KL}[2]{D_{\mathrm{KL}}\left(#1\;\midd\;#2\right)}

\def\1{\mathbf{1}}
\def\supp{\mathrm{supp}}

\def\eps{\varepsilon}

\def\opt{\mathrm{opt}}

\def\K{{\mathcal{K}}}
\def\P{{\mathrm{P}}}
\def\T{{\mathcal{T}}}
\def\X{{\mathcal{X}}}
\def\tokens{{\Sigma}}

\def\c{\mathbf{c}}
\def\r{\mathbf{r}}
\def\x{\mathbf{x}}
\def\endtoken{{\oslash}}

\newcommand{\adam}[1]{\textcolor{blue}{\textit{(#1)}}}
\newcommand{\david}[1]{\textcolor{red}{\textit{(#1)}}}

